\section{\fpfourres{} Analysis}
\label{sec:analysis}

Whether \fpfourres{} merits implementation depends on its post-training accuracy across model types. Table~\ref{tab:accuracy} reports the four implemented formats, NVFP4, and two single-sided ablations across nine models in three modalities.

\textbf{Quantisation Protocol.} Quantisation is post-training and calibration-free: weights use static per-block scaling ($k=32$, OCP), activations use dynamic scaling from each tensor's peak magnitude, and all formats round to nearest-even. Language model evaluations average two seeds; vision and speech use a single seed (Table~\ref{tab:accuracy}, footnote).

\begin{table*}[t]
  \centering
  \caption{Post-training quantisation accuracy. WikiText-2 perplexity
           ($\downarrow$), ImageNet-1k top-1 ($\uparrow$), LibriSpeech
           test-clean WER ($\downarrow$). All layers quantised; BF16 is the
           unquantised reference. The two rules below \fpfourres{} are
           single-sided and symmetric ablations of the residue construction.}
  \label{tab:accuracy}
  \setlength{\tabcolsep}{4pt}
  \renewcommand{\arraystretch}{1.1}
  \small
  \begin{tabular}{|l|r|r|r|r|r|r|r|}
    \hline
    & \multicolumn{4}{c|}{Perplexity $\downarrow$}
    & \multicolumn{2}{c|}{Top-1 $\uparrow$}
    & WER $\downarrow$ \\
    \hline
    Format & GPT-2 & Llama-2 & Llama-3 & Mistral & ViT-B\textdagger & RN-50\textdagger & wav2vec2\textdagger \\
           &       & 7\,B    & 8\,B    & 7\,B    &       &       &          \\
    \hline
    BF16   & 30.3 & 6.02 & 6.74 & 6.14 & 80.3 & 80.4 & 2.86 \\
    \hline
    \fpfour             & 107.8 & 7.83 & 9.19 & 7.90 & 73.6 & 11.3 & 9.00 \\
    \hline
    NVFP4               &  36.3 & \underline{6.38} & \underline{7.59} & 6.46 & 77.3 & 55.5 & \underline{3.68} \\
    \hline
    \mtwo               &  46.4 & 6.46 & 7.85 & \underline{6.35} & \underline{78.4} & 60.3 & 5.00 \\
    \hline
    \rowcolor{gray!12}
    \textbf{\fpfourres{} (residue on A)} & \textbf{\underline{35.4}} & \textbf{6.65} & \textbf{7.75} & \textbf{6.49} & \textbf{78.1} & \textbf{27.2} & \textbf{30.5} \\
    \hline
    \quad \textbf{\textit{residue on W}}    & \textbf{79.7} & \textbf{6.73} & \textbf{7.82} & \textbf{6.77} & \textbf{77.2} & \textbf{\underline{67.3}} & \textbf{4.27} \\
    \hline
    \quad \textit{residue on both} &  30.3 & 6.04 & 6.79 & 6.15 & 80.3 & 79.3 & 2.77 \\
    \hline
    \fpeight-E4M3       &  31.1 & 6.08 & 6.87 & 6.20 & 79.5 & 77.9 & 2.71 \\
    \hline
  \end{tabular}
  \\[2pt]
  \begin{minipage}{\textwidth}
    \footnotesize
    \underline{Underline}: best among MXFP4 / NVFP4 / \mtwo{} / \fpfourres{} (residue on A) per column (BF16 and \fpeight-E4M3 are reference points, not candidates).
    \fpfourres{} (default) places residue on activations only (W/A $=4.25/8.50$\,b); ``residue on W'' is mirrored ($8.50/4.25$\,b) and ``residue on both'' is symmetric ($8.50/8.50$\,b).
    \textdagger vision/speech: single seed (42).
  \end{minipage}
\end{table*}

Symmetric residue matches BF16 and \fpeight{} accuracy everywhere, confirming that two independent \fpfour{} blocks are an adequate substitute for one \fpeight{} block. Restricting the residue to activations alone, however, is not universally optimal: it recovers most of the accuracy gap on seven of nine evaluated models, including GPT-2, all three larger language models, and the vision transformer, but the preference inverts on the deeper convolutional network and the speech model, where naive \fpfour{} already collapses badly and weight-side residue instead dominates recovery. The inversion is not explained by architecture family or parameter volume: two closely related convolutional networks land on opposite sides despite near-identical weight-to-activation ratios. A per-layer error diagnostic offers real mechanistic insight (below), though it explains why an allocation fails more reliably than it predicts which allocation will win.

For the speech model, neither candidate explanation holds up: a sign-correlation test rules out error cancellation between weight and activation terms, and the error is not concentrated in the audio front end. The actual mechanism is subtler: activation-side residue collapses a large but symmetric, self-averaging quantisation error at one late, decoder-critical layer, but in doing so leaves a smaller, systematic per-channel weight bias as the dominant remaining error there instead. That structured bias proves more damaging to sequence decoding than the larger, unstructured noise it replaces. The result generalises beyond this format: reducing reconstruction error is not always sufficient when the reduction changes the character of the error, not just its size, at a decision-critical layer.

\fpfourres{} offers two single-sided configurations, placing the residue block on activations or on weights, and their strengths fall on different workloads. The activation-side variant achieves the lowest perplexity among all four-bit candidates on the compact decoder (GPT-2) and outperforms \mtwo{} on Llama-3-8B, though NVFP4 leads on the larger language models overall. The weight-side variant wins outright on the deeper convolutional network and outperforms \mtwo{} on the speech model. The two formats are therefore complementary rather than competing, motivating exposing both as independent coprocessor engines (\S\ref{sec:arch-engines}).

\subsection{Zero-shot task validation}
\label{sec:analysis-zeroshot}

\begin{table}[t]
  \centering
  \caption{Zero-shot task accuracy (\SI{}{\percent}) on downstream benchmark evaluations for GPT-2 and Llama-2-7B. Bold indicates the higher-performing single-sided residue allocation.}
  \label{tab:zeroshot}
  \setlength{\tabcolsep}{2.5pt}
  \renewcommand{\arraystretch}{1.1}
  \scriptsize
  \begin{tabular}{|l|l|r|r|r|r|r|}
    \hline
    Model & Task & BF16 & \fpfour & \mtwo & \fpfourres{} (Act) & \fpfourres{} (Wt) \\
    \hline
    GPT-2 & WinoGrande & 51.6 & 48.5 & 49.8 & \textbf{52.6} & 51.7 \\
    \hline
          & PIQA       & 62.3 & 58.2 & 60.1 & \textbf{61.4} & 58.3 \\
    \hline
          & HellaSwag  & 31.3 & 29.0 & 30.6 & \textbf{29.9} & 29.8 \\
    \hline
          & ARC-Easy   & 40.2 & 38.3 & 37.8 & \textbf{39.5} & 38.7 \\
    \hline
          & BoolQ\textdagger & 50.2 & 59.6 & 57.6 & 57.6 & 56.7 \\
    \hline
    Llama-2-7B & WinoGrande & 69.1 & 66.5 & 67.6 & \textbf{68.8} & 68.1 \\
    \hline
               & PIQA       & 78.4 & 76.9 & 77.8 & \textbf{77.9} & 77.5 \\
    \hline
               & HellaSwag  & 76.1 & 70.6 & 74.5 & \textbf{74.1} & 74.0 \\
    \hline
               & ARC-Easy   & 74.1 & 69.4 & 73.4 & \textbf{72.3} & 71.3 \\
    \hline
               & BoolQ      & 79.4 & 72.8 & 77.3 & \textbf{77.1} & 76.6 \\
    \hline
  \end{tabular}
  \\[2pt]
  \begin{minipage}{\columnwidth}
    \footnotesize
    \textdagger GPT-2 unquantised (BF16) accuracy on BoolQ is statistically indistinguishable from the \SI{50}{\percent} chance baseline; nominal gains under quantisation represent noise and are excluded from average gap recovery calculations.
  \end{minipage}
\end{table}


Five zero-shot reasoning benchmarks (Table~\ref{tab:zeroshot}) corroborate the perplexity-level findings: activation-side residue outperforms weight-side residue on four of five tasks for both the compact decoder and the instruction-tuned model, recovering \SI{78}{\percent} and \SI{67}{\percent} of the quantisation gap respectively (the compact decoder's near-chance BoolQ baseline excluded). \mtwo{} follows the same ordering under zero-shot evaluation as under perplexity, trailing activation-only residue on the compact decoder but competitive on the instruction-tuned model.